% Application Note for Bioinformatics (Oxford Academic).
%
% Target: 4 printed pages, Applications Note category.
%
% Oxford supplies a class file with its author kit; this document deliberately
% uses stock `article` so it compiles anywhere, and mirrors the Applications
% Note section order (Abstract / 1 Introduction / 2 Methods / 3 Results /
% 4 Discussion / Availability). Swap \documentclass for the journal class at
% submission time; no body text should need to change.
%
% Build:  latexmk -pdf docs/APPLICATION_NOTE.tex
%
% -------------------------------------------------------------------------
% EVERY NUMBER IN THIS PAPER COMES FROM A MEASURED RUN.
% The macros below are the single source of truth; they are transcribed from
% results/summary.txt and results/accuracy.txt, produced by
% scripts/benchmark.sh. Re-run the benchmark and update this block -- never
% edit a figure inline in the body text.
% -------------------------------------------------------------------------

\documentclass[11pt,a4paper]{article}

\usepackage[T1]{fontenc}
\usepackage[utf8]{inputenc}
\usepackage[margin=2.2cm]{geometry}
\usepackage{booktabs}
\usepackage{graphicx}
\usepackage{amsmath}
\usepackage{url}
\usepackage[hidelinks]{hyperref}
\usepackage{caption}
\usepackage{natbib}

% Deliberately core-LaTeX only: this compiles on a stock BasicTeX/TeX Live
% minimal install with no package downloads, which matters when the build host
% has no outbound network. siunitx and authblk are not assumed.

% --- Measured results (see header note) -----------------------------------
\newcommand{\NumRegions}{100{,}000}
\newcommand{\NumFragments}{4}            % million fragments
\newcommand{\NumAlignments}{8}           % million alignments
\newcommand{\BamSize}{87}                % MB
\newcommand{\NumCores}{12}
\newcommand{\WindowSpec}{$\pm$2\,kb}
\newcommand{\BinSize}{50}

\newcommand{\OursTime}{1.42}
\newcommand{\OursMem}{84.7}
\newcommand{\CovTime}{7.11}
\newcommand{\CovMem}{320.7}
\newcommand{\MatTime}{43.06}
\newcommand{\MatMem}{432.4}
\newcommand{\TotalTime}{50.17}
\newcommand{\SpeedupTotal}{35.4}
\newcommand{\SpeedupMatrix}{30.4}
\newcommand{\MemRatio}{5.1}

\newcommand{\ProfileR}{0.99999}
\newcommand{\MaxResidual}{0.38\%}
\newcommand{\MedianRowR}{0.988}
\newcommand{\FracAbove}{99.8\%}


\title{\vspace{-1.5cm}\textbf{cuttag\_profiler: air-gapped CUT\&Tag signal
profiling with an embedded browser dashboard}}

\author{Author One \and Author Two \\[2pt]
\small Department, Institution, City, Country}

\date{}

\begin{document}
\maketitle
\vspace{-1cm}

\begin{abstract}
\noindent\textbf{Summary:} Chromatin profiling assays such as CUT\&Tag and
CUT\&RUN are read out by quantifying normalised coverage over large sets of
genomic loci. The established tool for this step, \texttt{deeptools}, requires
a two-stage pipeline and a command-line interface, and takes tens of minutes on
a genome-scale locus set. Interactive alternatives are typically web services,
which many clinical and industrial laboratories cannot use because aligned
sequence data may not leave the premises. We present \texttt{cuttag\_profiler},
a single statically linked C\texttt{++}20 binary that computes reference-point
signal matrices directly from an indexed BAM and serves an interactive
dashboard from an embedded HTTP server bound to the loopback interface. On
\NumRegions{} promoter windows it is \SpeedupTotal$\times$ faster than the
equivalent \texttt{deeptools} pipeline while using \MemRatio$\times$ less
memory, and reproduces its meta-profile with a Pearson correlation of
\ProfileR. No byte of sequence data crosses a network boundary.

\smallskip
\noindent\textbf{Availability and implementation:} Source, build instructions
and the benchmark harness are available at
\url{https://github.com/ORGANISATION/cuttag_profiler} under the MIT licence.
The binary depends only on \texttt{htslib} and the platform C library.

\smallskip
\noindent\textbf{Contact:} \url{author@institution.edu}
\end{abstract}

\section{Introduction}

CUT\&Tag \citep{kayajaya2019} and CUT\&RUN \citep{skene2017} localise a histone
modification or transcription factor by tethering an enzyme to a target-bound
antibody. Downstream interpretation almost always reduces to the same
operation: for each of $N$ loci, compute normalised read coverage in $B$ bins
across a fixed window anchored on a reference point such as a transcription
start site, then summarise the resulting $N \times B$ matrix as a meta-profile
and a heatmap.

Two practical problems attach to that operation.

First, cost. The \texttt{deeptools} suite \citep{ramirez2016} performs it in two
stages: \texttt{bamCoverage} converts the BAM into a bigWig coverage track, and
\texttt{computeMatrix} windows that track. Each stage rewrites the data, and the
intermediate track is written at genome scale even when the loci of interest
cover a small fraction of the genome.

Second, access. The scientists who interpret these matrices are usually not the
scientists who can operate a command line, and the interactive tools that close
that gap are predominantly hosted services. Laboratories operating under
GDPR, HIPAA or an institutional data-governance policy frequently prohibit
transmitting aligned human sequence data to third-party infrastructure, which
disqualifies hosted viewers regardless of their merits.

\texttt{cuttag\_profiler} addresses both. It computes the matrix in one pass
directly from the indexed BAM, and it ships its own dashboard inside the
executable, served over the loopback interface only.

\section{Methods}

\subsection{Architecture}

The tool is a single C\texttt{++}20 executable with no runtime dependency
beyond \texttt{htslib} and the platform C library; \texttt{htslib} is linked
statically, so deployment is a file copy. It has two modes: a headless
\texttt{profile} subcommand for pipelines, and a \texttt{serve} subcommand that
starts the dashboard.

\subsection{Parallel indexed query}

Work is distributed one locus per task over a pool of worker threads, block
partitioned so that each worker writes a contiguous stripe of the output matrix.
\texttt{htslib}'s index and header are immutable once loaded and are shared;
its \texttt{htsFile} handles carry mutable decompression state and are not, so
each worker leases a private handle from a pool.

A design decision worth recording is what \emph{not} to share. \texttt{htslib}
offers a thread pool that parallelises BGZF block inflation, and the intuitive
configuration is one pool sized to the machine. For this access pattern that is
actively harmful: with per-region parallelism already saturating the cores,
routing every block through a shared pool serialises inflation behind the
pool's queue and adds a hand-off per block. Measured on the benchmark below, a
\NumCores-thread shared pool cost 5.95\,s wall time against
1.38\,s with no pool at all, and consumed 52\,s of CPU
against 6.4\,s. The default is therefore no shared pool: each
handle inflates on the thread that owns it. A pool is created only when the
caller explicitly requests single-threaded operation, where it can still
overlap inflation within one query.

Per-thread reduction buffers are allocated on 64\,byte boundaries and
padded to whole cache lines, so no two workers share a line while accumulating
column sums.

\subsection{Signal quantification}

For a locus with anchor $a$, the window $[a - u, a + d)$ is divided into
$B = \lceil (u+d)/w \rceil$ bins of width $w$. For minus-strand features the
window is mirrored and the resulting row reversed, so that bin index always
runs $5' \rightarrow 3'$.

Each alignment contributes to the bins its reference span overlaps. For
properly paired data the span is the full sequenced fragment, inferred from
\texttt{TLEN} and counted once through the leftmost mate; this is the correct
footprint for CUT\&Tag, where the fragment marks the protected region. Because
a fragment may begin before the window and still intersect it, the underlying
iterator query is widened to the left by the maximum fragment length; omitting
this silently truncates signal at window edges.

Two tallies are accumulated per bin: the number of overlapping fragments, and
the number of aligned bases, the latter honouring the CIGAR string so that
deletions and reference skips do not inflate coverage. Reported signal is
either fragments per bin (the default, matching \texttt{deeptools}) or mean
per-base depth. Normalisation is one of raw, CPM, RPKM or BPM; the library
size used as the denominator is read from the index metadata in
$O(\text{contigs})$ rather than by streaming the file.

\subsection{Embedded dashboard}

\texttt{serve} starts an HTTP server (\texttt{cpp-httplib}) bound to
\texttt{127.0.0.1} and opens the user's browser. Three properties enforce the
air gap:

\begin{enumerate}\itemsep2pt
  \item the listener refuses to bind a non-loopback address unless the operator
        passes an explicit override flag, which also prints a warning;
  \item every asset is served from a byte array compiled into the binary by a
        CMake generator step, so no CDN, font host or analytics endpoint is
        contacted;
  \item responses carry \texttt{Content-Security-Policy: default-src 'self'},
        so even a modified asset has no route out.
\end{enumerate}

The dashboard posts a JSON query description to \texttt{/api/profile} and
receives the meta-profile, a band-averaged heatmap and run statistics. Full
matrices are retrieved separately as TSV, keyed by a run token, rather than
serialised into every interactive response.

\section{Results}

\subsection{Benchmark design}

We compared against \texttt{deeptools} \texttt{computeMatrix reference-point}.
Because \texttt{computeMatrix} consumes a bigWig rather than a BAM, the
comparable unit of work is the two-stage pipeline an analyst actually runs:
\texttt{bamCoverage} followed by \texttt{computeMatrix}. Both stages are timed,
and the speedup is quoted against their sum; the \texttt{computeMatrix} stage is
also reported alone, which is the appropriate comparison for a reader who
already holds a cached coverage track.

The dataset is a synthetic paired-end CUT\&Tag library of \NumFragments{}
million fragments (\NumAlignments{} million alignments, \BamSize{}~MB BAM) over
a 250\,Mb chromosome, with \NumRegions{} annotated promoters, of
which 55\% of fragments are drawn from a TSS-centred peak
($\sigma = 180\,bp$) and the remainder from uniform background. It is
generated by \texttt{scripts/make\_demo\_data.py} at a fixed seed, so the
benchmark is reproducible without redistributing controlled-access data. Both
tools were given \NumCores{} threads and \WindowSpec{} windows in
\BinSize\,bp bins, three replicates each, on an Apple M-series
workstation.

\subsection{Runtime and memory}

\begin{table}[htbp]
\centering
\caption{Wall-clock time and peak resident memory over \NumRegions{} promoter
windows, \NumCores{} threads, mean of three replicates.}
\label{tab:benchmark}
\begin{tabular}{lrr}
\toprule
Stage & Time (s) & Peak RSS (MB) \\
\midrule
\texttt{cuttag\_profiler} (BAM $\rightarrow$ matrix) & \OursTime & \OursMem \\
\midrule
\texttt{deeptools bamCoverage} (BAM $\rightarrow$ bigWig) & \CovTime & \CovMem \\
\texttt{deeptools computeMatrix} (bigWig $\rightarrow$ matrix) & \MatTime & \MatMem \\
\textit{deeptools pipeline total} & \textit{\TotalTime} & \textit{\MatMem} \\
\bottomrule
\end{tabular}
\end{table}

Table~\ref{tab:benchmark} gives \SpeedupTotal$\times$ end-to-end and
\SpeedupMatrix$\times$ against \texttt{computeMatrix} alone, at
\MemRatio$\times$ lower peak memory. The memory difference follows from the
architecture: no genome-wide intermediate track is materialised, so resident
memory scales with the requested matrix rather than with the genome.

\subsection{Concordance with \texttt{deeptools}}

Agreement was assessed on the same matrices
(\texttt{scripts/compare\_matrices.py}). The two tools normalise against
different denominators --- \texttt{deeptools} scales binned coverage derived
from a bigWig, \texttt{cuttag\_profiler} scales fragment counts taken directly
from the BAM --- so the tools are interchangeable if and only if the signal
\emph{shape} matches up to one global scale factor, which is what we test.

The meta-profiles correlate at $r = \ProfileR$, with a maximum residual of
\MaxResidual{} of peak signal after fitting a single global scale factor.
Per-locus agreement, over a systematic sample of matched rows, has median
$r = \MedianRowR$, and \FracAbove{} of loci exceed $r = 0.95$. Residual
disagreement is attributable to the intermediate binning \texttt{deeptools}
performs when building the bigWig, which the direct query avoids.

\section{Discussion}

The speedup has a structural rather than a micro-optimisation explanation.
\texttt{deeptools} materialises a genome-wide coverage track and then windows
it, so its cost scales with genome size; \texttt{cuttag\_profiler} queries only
the intervals requested, so its cost scales with the loci of interest. Whenever
those loci cover a modest fraction of the genome --- the usual case for
promoter, enhancer or peak-set analyses --- the direct query does
proportionally less work. The gap narrows when the locus set approaches whole
genome coverage, and a cached bigWig that is reused across many locus sets
amortises the conversion, which is why Table~\ref{tab:benchmark} reports the
stages separately.

Two limitations should be stated. The benchmark uses simulated data, chosen so
the comparison is reproducible without redistributing controlled-access
alignments; the concordance analysis, not an absolute enrichment value, is what
the simulation is asked to support. Second, BPM normalisation is computed over
the supplied locus set rather than genome-wide, since no genome-wide pass is
performed; where a genome-wide BPM denominator is required, CPM or RPKM should
be used instead.

The design generalises past CUT\&Tag: nothing in the query or scoring layer is
assay-specific, and the same binary profiles ChIP-seq or ATAC-seq alignments.
The air-gap property is the part we would emphasise for adoption. A tool that
is fast but requires a data-transfer agreement is, in a regulated laboratory,
not usable at all; one that runs entirely on the analyst's own workstation
needs no such agreement.

\section*{Acknowledgements}

We thank the \texttt{htslib} and \texttt{deeptools} developers, whose work this
tool builds on and measures itself against.

\section*{Funding}

This work was supported by [funder, grant number].

\begin{thebibliography}{9}

\bibitem[Kaya-Okur \textit{et al.}(2019)]{kayajaya2019}
Kaya-Okur,~H.S. \textit{et al.} (2019)
CUT\&Tag for efficient epigenomic profiling of small samples and single cells.
\textit{Nat. Commun.}, \textbf{10}, 1930.

\bibitem[Skene and Henikoff(2017)]{skene2017}
Skene,~P.J. and Henikoff,~S. (2017)
An efficient targeted nuclease strategy for high-resolution mapping of DNA
binding sites. \textit{eLife}, \textbf{6}, e21856.

\bibitem[Ram{\'i}rez \textit{et al.}(2016)]{ramirez2016}
Ram{\'i}rez,~F. \textit{et al.} (2016)
deepTools2: a next generation web server for deep-sequencing data analysis.
\textit{Nucleic Acids Res.}, \textbf{44}, W160--W165.

\bibitem[Bonfield \textit{et al.}(2021)]{bonfield2021}
Bonfield,~J.K. \textit{et al.} (2021)
HTSlib: C library for reading/writing high-throughput sequencing data.
\textit{GigaScience}, \textbf{10}, giab007.

\bibitem[Li \textit{et al.}(2009)]{li2009}
Li,~H. \textit{et al.} (2009)
The Sequence Alignment/Map format and SAMtools.
\textit{Bioinformatics}, \textbf{25}, 2078--2079.

\end{thebibliography}

\end{document}
